\section{Questão 15}

Considere a cadeia de Markov com três estados $S = \{1, 2, 3\}$, 
que tem o diagrama de transição de estado mostrado na figura a seguir
%
\begin{figure}[htbp]
    \centering
    \begin{tikzpicture}
        \node[state] (1) {1};
        \node[state, right=3 of 1] (2) {2};
        \node[state, below right=1 and 1.4 of 1] (3) {3};

        \path [->]
            (1) edge[loop above] node [above left] {$1/4$} (1)
            (1) edge[bend left] node [above] {$3/4$} (3)
            %
            (2) edge[bend right] node [above] {$1/2$} (1)
            (2) edge[bend left] node [below right] {$1/2$} (3)
            %
            (3) edge[bend left] node [below left] {$1/2$} (1)
            (3) edge[bend left] node [above] {$1/4$} (2)
            (3) edge[loop below] node [below] {$1/4$} (3);
    \end{tikzpicture}
\end{figure}

Suponha que $P(X_0=1) = 1/2$ e $P(X_0=2) = 1/4$. 
%
\begin{itemize}
    \item [a)] Encontre a matriz de transição de estado para esta cadeia de Markov.
    \item [b)] Determine $P(X_0=3, X_1=2, X_2=1)$.
    \item [c)] Encontre $P(X_0=3,X_1=1, X_2=1, X_3=3)$.
\end{itemize}

\begin{respostas}
\begin{itemize}
    \item [a)] Por inspeção, conclui-se que a matriz de transição de probabilidade é dada por
    %
    \begin{equation}
        P = \begin{bmatrix}
            1/4 & 0   & 3/4 \\
            1/2 & 0   & 1/2 \\
            1/2 & 1/4 & 1/4
        \end{bmatrix}.
    \end{equation}

    \item [b)] Tem-se
    %
    \begin{align}
        P(X_0=3, X_1=2, X_2=1) &= P(X_0=3) \cdot P(X_1=2 | X_0=3) \cdot P(X_2=1 | X_1=2 ) \\
        &= \left(1 - \frac{1}{2} - \frac{1}{4}\right) \cdot \frac{1}{4} \cdot \frac{1}{2} \\
        &= \frac{1}{32}.
    \end{align}
    
    \item [c)] Tem-se
    %
    \begin{align}
        P &= P(X_0=3) \cdot P(X_1=1 | X_0=3) \cdot P(X_2=1 | X_1=1) \cdot P(X_3=3 | X_2=1) \\
        &= \left(1 - \frac{1}{2} - \frac{1}{4}\right) \cdot \frac{1}{2} \cdot \frac{1}{4} \cdot \frac{3}{4} \\
        &= \frac{3}{128}.
    \end{align}
\end{itemize}
\end{respostas}